\documentclass[12pt]{article}
\usepackage[utf8]{inputenc}
\usepackage[T1]{fontenc}
\usepackage{hyperref}
\title{Tai Long Language Resources for K-12 Curriculum Development}
\author{Compiled by Claude Sonnet 5 (claude-sonnet-5) for Andr\'as Kornai}
\date{}
\begin{document}
\maketitle

\section{Introduction}
Tai Long (ISO 639-3: thi) is a Kra-Dai language of northern Laos. The search found almost nothing about the language itself. The ethnonym ``Tai Long'' overlaps with Shan names, and the sources found do not settle how the two relate, so identification should be checked against Ethnologue before any use.

\section{Primary Educational Resources}
None found.

\section{Linguistic and Reference Materials}
None specific to Tai Long were found. The Chulalongkorn Center of Excellence in Southeast Asian Linguistics lists Tai Long (Shan) among Southwestern Tai languages:
\url{https://www.arts.chula.ac.th/chulaseal/publication/tai-languages/}

\section{Religious and Community Resources}
None found in the search results.

\section{Digital and Audio Resources}
None found in the search results.

\section{Curriculum Development Context}
Joshua Project reports about 2,300 Tai Long in Louangphabang Province, citing Ethnologue 2016, with slight dialect differences between locales:
\url{https://joshuaproject.net/people_groups/18690}
No evidence of school use was found. Lao is the national language and school language in Laos.
\end{document}